% !PRINTED 382
\weberpage{382}
%<<CONT
whose roots are the ten roots syzygetic with $\xi_1\xi_2$. If $\xi_3$ denotes one of these, then

\weq{3}{X(\xi_1,\xi_2,\xi_3) = 0,}

and from this follows the theorem

\gesperrt{1. There is a rational function of three variables $X(\xi_1,\xi_2,\xi_3)$ which vanishes when $\xi_1\xi_2\xi_3$ is any syzygetic triple of roots of the double tangent equation, and which does not vanish when $\xi_1\xi_2\xi_3$ is an azygetic triple.}

By \S.\,100, by means of a complete seven-system of these roots

\weq{4}{\xi_1,\ \xi_2,\ \xi_3,\ \xi_4,\ \xi_5,\ \xi_6,\ \xi_7}

all the roots can be expressed rationally, and indeed in such a way that, e.g.,

\weq{5}{\xi_{12} = \Psi(\xi_1,\xi_2 \mid \xi_3,\xi_4,\xi_5,\xi_6,\xi_7)}

becomes a root, where $\Psi$ denotes a rational function which does not change when $\xi_1$ and $\xi_2$ are interchanged, or when $\xi_3,\xi_4,\xi_5,\xi_6,\xi_7$ are permuted arbitrarily. But if in (5), the function $\Psi$ being held fixed, another pair $\xi_i,\xi_k$ is put in place of $\xi_1,\xi_2$, one obtains another root $\xi_{ik}$. The roots (4) may also be denoted by $[1\,8],\dots[7\,8]$ and $\xi_{ik}$ by $[i\,k]$.

By every permutation of the roots (1), according to (2), the whole system of the roots $[i\,k]$ undergoes a certain permutation.

But if one replaces the complete seven-system (4) by another, the formula (5) yields a definite other root, and the whole system of the roots $[i\,k]$ is subjected to a second permutation. It is then first of all easy to prove:

\gesperrt{2. The permutation group $P$ of the roots of the double tangent equation, which one obtains when in (4) and (5) one puts in place of $\xi_1,\xi_2,\dots\xi_7$ all complete systems, each in every arbitrary order, is the Galois group of the double tangent equation.}

To demonstrate this, we have two things to show:
%CONT>>
